\documentclass[10pt,a4paper]{amsart}
\usepackage{amsmath,amssymb,amsthm,graphicx}
\usepackage{verbatim}
\newenvironment{reference}{\comment}{\endcomment}
%\newenvironment{reference}{}{}
\newenvironment{slogan}{\comment}{\endcomment}
\newenvironment{history}{\comment}{\endcomment}

% For commutative diagrams we use Xy-pic
\usepackage[all]{xy}

% We use 2cell for 2-commutative diagrams.
\xyoption{2cell}
% Theorem environments.
%
\theoremstyle{plain}
\newtheorem{theorem}[subsection]{Theorem}
\newtheorem{proposition}[subsection]{Proposition}
\newtheorem{lemma}[subsection]{Lemma}

\theoremstyle{definition}
\newtheorem{definition}[subsection]{Definition}
\newtheorem{example}[subsection]{Example}
\newtheorem{exercise}[subsection]{Exercise}
\newtheorem{situation}[subsection]{Situation}

\theoremstyle{remark}
\newtheorem{remark}[subsection]{Remark}
\newtheorem{remarks}[subsection]{Remarks}

\numberwithin{equation}{subsection}

% Macros
%
\def\lim{\mathop{\mathrm{lim}}\nolimits}
\def\colim{\mathop{\mathrm{colim}}\nolimits}
\def\Spec{\mathop{\mathrm{Spec}}}
\def\Hom{\mathop{\mathrm{Hom}}\nolimits}
\def\Ext{\mathop{\mathrm{Ext}}\nolimits}
\def\SheafHom{\mathop{\mathcal{H}\!\mathit{om}}\nolimits}
\def\SheafExt{\mathop{\mathcal{E}\!\mathit{xt}}\nolimits}
\def\Sch{\mathit{Sch}}
\def\Mor{\mathop{\mathrm{Mor}}\nolimits}
\def\Ob{\mathop{\mathrm{Ob}}\nolimits}
\def\Sh{\mathop{\mathit{Sh}}\nolimits}
\def\NL{\mathop{N\!L}\nolimits}
\def\CH{\mathop{\mathrm{CH}}\nolimits}
\def\proetale{{pro\text{-}\acute{e}tale}}
\def\etale{{\acute{e}tale}}
\def\QCoh{\mathit{QCoh}}
\def\Ker{\mathop{\mathrm{Ker}}}
\def\Im{\mathop{\mathrm{Im}}}
\def\Coker{\mathop{\mathrm{Coker}}}
\def\Coim{\mathop{\mathrm{Coim}}}

% Boxtimes
%
\DeclareMathSymbol{\boxtimes}{\mathbin}{AMSa}{"02}

%
% Macros for moduli stacks/spaces
%
\def\QCohstack{\mathcal{QC}\!\mathit{oh}}
\def\Cohstack{\mathcal{C}\!\mathit{oh}}
\def\Spacesstack{\mathcal{S}\!\mathit{paces}}
\def\Quotfunctor{\mathrm{Quot}}
\def\Hilbfunctor{\mathrm{Hilb}}
\def\Curvesstack{\mathcal{C}\!\mathit{urves}}
\def\Polarizedstack{\mathcal{P}\!\mathit{olarized}}
\def\Complexesstack{\mathcal{C}\!\mathit{omplexes}}
% \Pic is the operator that assigns to X its picard group, usage \Pic(X)
% \Picardstack_{X/B} denotes the Picard stack of X over B
% \Picardfunctor_{X/B} denotes the Picard functor of X over B
\def\Pic{\mathop{\mathrm{Pic}}\nolimits}
\def\Picardstack{\mathcal{P}\!\mathit{ic}}
\def\Picardfunctor{\mathrm{Pic}}
\def\Deformationcategory{\mathcal{D}\!\mathit{ef}}
\usepackage{cleveref}
\crefname{originalreference}{original reference}{original references}
\Crefname{originalreference}{Original reference}{Original references}
\makeatletter
\newlabel{lemma-fully-faithful-spaces-over-X}{{096K}{source}{}{sourcepointer-1}{}}
\newlabel{lemma-fully-faithful-spaces-over-X@cref}{{[originalreference][0][]096K}{[1][0][]source}{}{}{}}
\makeatother
\crefname{originalreference}{original reference}{original references}
\Crefname{originalreference}{Original reference}{Original references}
\makeatletter
\newlabel{lemma-make-w-local}{{0975}{source}{}{sourcepointer-2}{}}
\newlabel{lemma-make-w-local@cref}{{[originalreference][0][]0975}{[1][0][]source}{}{}{}}
\makeatother
\title{Universal property of the w-localization of a ring}
\author{The Stacks Project Authors}
\date{stacks-project:proetale.tex:local-snapshot-20261012; GNU FDL 1.2 or later}
\begin{document}
\maketitle
\noindent\textit{Editorial note.} All delivered mathematical text is the author's own, copied from the pinned Stacks Project chapter file proetale.tex; no statement or proof step was rewritten, reordered or removed.\par\medskip
\noindent\textit{Editorial note.} The counted proof is the author's complete argument: the reduction to a factorization through a partial w-localization, the decomposition of the target into open and closed pieces through the constructible strata and the profinite structure of the closed point set, and the conversion of the unique continuous map into a unique morphism of schemes.\par\medskip
\noindent\textit{Editorial note.} The label lemma-universal names the counted statement, and the proof environment immediately following it in the source is delivered verbatim as its proof.\par\medskip
\noindent\textit{Editorial note.} The counted statement is carried by the author's own lemma environment in the chapter source.\par\medskip
\noindent\textit{Editorial note.} Pre-open gate: python operations/agent-authoring/tools/precard\_gate.py --source proetale.tex --label lemma-universal --start 913 --end 946 --local-id 3589 returned RESULT=OK (L3589 free, occupancy FREE, public interval clear). Candidate taken from sources/stacks-proetale/candidates-halfA.json.\par\medskip
\noindent\textit{Editorial note.} The statement and proof are copied from the pinned chapter file; the cited inputs are kept as external pointers with their canonical Stacks tags.\par\medskip
\noindent\textit{Allowed standard prerequisites (editorial).}\begin{itemize}
\item Structure of the partial w-localizations. Same chapter, the construction of the w-localization as the colimit of the A\_E and the description of the closed points as the limit of the closed subsets Z\_E, with the strata Z(E',E'') locally closed and constructible. Supplies the reduction to a unique factorization through some A\_E mapping the closed points into Z\_E, and the description of the strata.
\item Profinite structure of the closed points of a w-local scheme. Same chapter, the lemma that the set of closed points of a w-local scheme is profinite, so that every constructible subset of it is open and closed. Supplies the decomposition of Y\_0 into the open and closed pieces f\textasciicircum{}\{-1\}(Z(E',E'')) cap Y\_0 and the identification of Y\_0 with the set of connected components of Y.
\item Faithfulness of the functor from schemes to spaces over X in the relevant case. Same chapter, the lemma that the map on underlying topological spaces determines the morphism over X for the schemes involved. Supplies the unique morphism of schemes Y -> X\_\{E',E''\} over X from the unique continuous map.
\end{itemize}\medskip
\ifdefined\hypertarget\hypertarget{sourcepointer-1}{}\fi
\noindent\textit{Source reference.} Original reference 096K: sources/stacks-proetale/pinned-original-source.tex, source lines 582--582. Prerequisite pointer invoked at proof line(s) 943; it names a standard fact of the same chapter that the counted argument uses, and every step of the counted proof built on it is retained verbatim in the delivered body.\par\medskip
\ifdefined\hypertarget\hypertarget{sourcepointer-2}{}\fi
\noindent\textit{Source reference.} Original reference 0975: sources/stacks-proetale/pinned-original-source.tex, source lines 844--844. Prerequisite pointer invoked at proof line(s) 916; it names a standard fact of the same chapter that the counted argument uses, and every step of the counted proof built on it is retained verbatim in the delivered body.\par\medskip
\begin{lemma}[Universal property of the construction]
\label{lemma-universal}
Let $A$ be a ring. Let $A \to A_w$ be the ring map constructed in
Lemma \ref{lemma-make-w-local}. For any ring map $A \to B$ such that
$\Spec(B)$ is w-local, there is a unique factorization $A \to A_w \to B$
such that $\Spec(B) \to \Spec(A_w)$ is w-local.
\end{lemma}

\begin{proof}
Denote $Y = \Spec(B)$ and $Y_0 \subset Y$ the set of closed points.
Denote $f : Y \to X$ the given morphism.
Recall that $Y_0$ is profinite, in particular every constructible
subset of $Y_0$ is open and closed. Let $E \subset A$ be a finite subset.
Recall that $A_w = \colim A_E$ and that the set of closed points of
$\Spec(A_w)$ is the limit of the closed subsets $Z_E \subset X_E = \Spec(A_E)$.
Thus it suffices to show there is a unique factorization $A \to A_E \to B$
such that $Y \to X_E$ maps $Y_0$ into $Z_E$.
Since $Z_E \to X = \Spec(A)$ is bijective, and since the strata
$Z(E', E'')$ are constructible we see that
$$
Y_0 = \coprod f^{-1}(Z(E', E'')) \cap Y_0
$$
is a disjoint union decomposition into open and closed subsets.
As $Y_0 = \pi_0(Y)$ we obtain a corresponding decomposition of
$Y$ into open and closed pieces. Thus it suffices to construct
the factorization in case $f(Y_0) \subset Z(E', E'')$ for
some decomposition $E = E' \amalg E''$.
In this case $f(Y)$ is contained in the set of points of $X$
specializing to $Z(E', E'')$ which is homeomorphic to $X_{E', E''}$.
Thus we obtain a unique continuous map $Y \to X_{E', E''}$ over $X$. By
Lemma \ref{lemma-fully-faithful-spaces-over-X}
this corresponds to a unique morphism of schemes
$Y \to X_{E', E''}$ over $X$. This finishes the proof.
\end{proof}

\end{document}
